
% ════════════════════════════════════════════
% 具体分析（5.2.1.具体分析.tex，三级子块）—— 写作提示词见本文件末尾【本章节写作提示词】。
% ════════════════════════════════════════════
\subsubsection{具体分析}

针对问题二，本问属于定位误差传播建模与多目标车辆调度优化类问题，采用线性误差传播解析建模、站对网格枚举与时序仿真组合优化相结合的方法进行求解。

首先由两条方位线的几何交会建立定位误差与基线长度、交会角、目标距离之间的关系，接着定义交会定位精度指标并在行驶时间约束下枚举告警 T2 的可观测站对，然后针对三告警并发场景建立车辆指派与时序仿真模型以最小化最大观测开始时间，最后与独立定义的最近指派启发式进行对比分析。

具体分析流程如图\ref{fig:analysis2_flow}所示。

\begin{figure}[ht]
    \begin{center}
        \begin{tikzpicture}[x=1cm, y=0.7cm]
            \node[box] (n11) at (0, 0) {交会误差建模};
            \node[box] (n12) at (5, 0) {精度指标定义};
            \node[box] (n13) at (10, 0) {可观测站枚举};
            \node[box] (n22) at (5, -3) {行驶时间约束};
            \node[box] (n21) at (0, -3) {最优站对求解};
            \node[box] (n23) at (10, -3) {车辆指派};
            \node[box] (n31) at (0, -6) {时序仿真};
            \node[box] (n32) at (5, -6) {最小化最大时间};
            \node[box] (n33) at (10, -6) {启发式对比};
            \draw[arrow] (n11) -- (n12);
            \draw[arrow] (n12) -- (n13);
            \draw[arrow] (n23) -- (n22);
            \draw[arrow] (n22) -- (n21);
            \draw[arrow] (n13) -- ++(0,-1.0) -| (n23);
            \draw[arrow] (n21) -- (n31);
            \draw[arrow] (n31) -- (n32);
            \draw[arrow] (n32) -- (n33);
        \end{tikzpicture}
        \caption{问题二分析流程图}
        \label{fig:analysis2_flow}
    \end{center}
\end{figure}
